\section{Introduction}
A prepared experiment determines which future tests can be executed, which histories have the same predictive content, and which distributions on that content are actually acquired. The resource problem is to turn those predictions into one causal program with finite persistent state. Formal parameter dimension is not enough: states may be reachable but scarcely acquired, a slow dynamical mode may be invisible to the task, and a good terminal quantizer need not admit an online implementation. This paper connects those distinctions through a two-sided physical-acquisition theorem.

Physical preparation need not erase the observer's evidence about an unknown parameter. With retained memory, successive physical excursions form a finite-boundary Markov-renewal process, not a sequence of independent cycles. The physical state and the reports may nevertheless be infinite. Several closed boundary classes and periodic behavior are allowed. Gain is obtained by forming a reward-to-duration ratio in each closed class and only then mixing by the actual absorption probabilities. Across a model family, the recurrent structure may change without damaging a specified task.

An approximate Poisson equation provides an upper estimate, but an upper estimate alone does not identify the resource obstruction. Let $P$ be the exit-label matrix, $T_{ij}=\E_i[\tau\one_{\{I_1=j\}}]$ the joint duration--exit matrix, $r$ the complete-excursion reward vector, and $g$ the physical-time gain. Put
\begin{equation}\label{eq:intro-poisson}
 b=r-Tg,\qquad
 k_N=\inf_u\left\{\norm{b-(I-P)u}_\infty+\frac{\osc(u)}N\right\}.
\end{equation}
The central result, Theorem~\ref{thm:modulus}, proves both that $k_N$ controls actual physical-acquisition discrepancy and that physical-acquisition discrepancies control $k_N$. The reverse comparison uses the largest prefix deviation, normalized by the total acquisition budget, and includes precisely an actual return-tail remainder. Periodic cancellation shows why the prefix cannot simply be replaced by the terminal time. A deterministic long excursion shows why boundary calculations alone cannot remove the physical-time remainder.

The converse mechanism is a finite renewal identity. Averaging the actual prefix deviations constructs a corrector whose residual is the sum of a remaining physical-tail reward and a duration-weighted terminal strip. It never separates exit labels from their durations. An exact linear-programming dual then gives signed approximate-invariance certificates for a lower bound on physical discrepancy. These certificates are not fictitious acquired distributions. Nor are they parameters or exact posterior coordinates provided to the observer. The result requires finite mean excursions uniformly over the declared entering rows, but no spectral gap, irreducibility or uniform group-inverse bound.

For executable predictive occupation tests, the same theorem becomes a two-sided temporal certificate for the actual acquired law. Theorem~\ref{thm:resolution} then localizes a small-ball witness to the intended spatial resolution $r$. A truncated squared-distance test transfers the mass whenever the occupation discrepancy is of order $r$, rather than requiring an error of order $r^2$. It retains the true mass of each lower-dimensional or singular witness. The resulting finite-acquisition lower bound combines with global candidate-valid covers and actual-law update stability to give the counted causal upper in Corollary~\ref{cor:resolution-memory}. Spatial recursive implementability remains a separately verified condition; the temporal converse does not conceal it.

Theorem~\ref{thm:profile} shows that the general modulus is attained against every history-dependent strategy on a continuum of primitive diagnostic experiments. Revelation intensity and task amplitude can be arbitrary continuous profiles. Two retained labels attain the exact minimax value at every horizon. Power profiles recover algebraic rates, while a logarithmic amplitude gives a nonalgebraic rate despite compactness and continuous gain. A queue with unbounded workload and a noncommuting quantum instrument give structurally different raw verifications. The latter crosses a preparation-rank change without a uniform recurrence bound. Neither realization is a premise of the general theorem.

For one completely specified scored experiment with fresh geometric marks, the same mechanisms give
\begin{equation}\label{eq:intro-law}
 \calR(4n,M,\delta,e)
 \asymp q_M(\rho)+\Psi_n(a,\lambda)+\delta^2
              +A_*(\varepsilon_*(e)),
\end{equation}
under explicit profile and quantization doubling. Here $\rho$ is the actually loaded mark law, $\Psi_n$ is the penalized marked modulus attained by the diagnostic task, $\delta$ is a fixed invisible calibration ambiguity, and $\varepsilon_*(e)$ is the causal deficiency of a real report-erasure interface. The exact upper and lower before doubling are also given. All four calls, controller labels, output labels and report deadlines are charged. Program description, temporary workspace and numerical precision are not identified with persistent labels.

The qualitative task-compatibility theorem is retained in Section~\ref{sec:transfer}: on a compact continuous-response family with uniformly integrable returns, continuity of gain is equivalent to uniform physical-acquisition and normalized-discount limits and uniformly finite approximate-corrector prices. Compact robust attainment and finite-model determination are consequences on a whole legal class satisfying that criterion. Finite semialgebraic instruments admit a separate effective criterion and recurrence-exhaustion schedule. Those statements do not imply unrestricted attainment when task compatibility fails.

\subsection*{Relation to earlier theory}
Multichain evaluation, classwise ratios and semi-Markov bias equations are classical~\cite{DenardoFox,FedergrunSchweitzer,LamondPuterman,Meyer}. They are not presented as new definitions of causal geometry. The literature on ergodic approximation explicitly relates convergence rates to residual--solution approximation functionals~\cite{Shaw}, and Poisson solutions generate martingale representations under established integrability conditions~\cite{GlynnInfanger}. Our contribution is the marked physical-call converse, its finite boundary-balance dual and its finite-resolution use under a fully specified causal resource model, rather than an assertion that approximation functionals or Poisson martingales are new.

Minimal predictive states and comparison of experiments likewise have substantial antecedents~\cite{ShaliziCrutchfield,Blackwell,Fritz}. The Borel realization hypotheses of our quotient lemma are explicit. Quantization and finite-memory learning are established subjects~\cite{GrafLuschgy,HellmanCover}; here actual occupation and the charged implementation are identified before these tools are used. Semialgebraic asymptotics and choice have important predecessors~\cite{BewleyKohlberg,BewleyKohlbergStationary,BasuPollackRoy}. Fixed-response linear programming and finite-algebraic certification are distinguished from efficient global synthesis. The separate primitive verification and counterexample arguments make these boundaries mathematical, rather than terminological.
